\documentclass[11pt]{article}
\usepackage[letterpaper,margin=0.72in]{geometry}
\usepackage[T1]{fontenc}
\usepackage[utf8]{inputenc}
\usepackage{lmodern}
\usepackage{microtype}
\usepackage{xcolor}
\usepackage{graphicx}
\usepackage{booktabs}
\usepackage{longtable}
\usepackage{tabularx}
\usepackage{array}
\usepackage{enumitem}
\usepackage[most]{tcolorbox}
\usepackage{fancyhdr}
\usepackage{titlesec}
\usepackage{hyperref}
\usepackage{xurl}
\usepackage{lastpage}
\usepackage{ragged2e}
\usepackage{multicol}
\usepackage{amssymb}
\definecolor{BAOrange}{HTML}{E87521}
\definecolor{BANavy}{HTML}{17324D}
\definecolor{BASlate}{HTML}{5B6573}
\definecolor{BALight}{HTML}{F4F6F8}
\definecolor{BAGold}{HTML}{D7A62A}
\hypersetup{colorlinks=true,linkcolor=BANavy,urlcolor=BAOrange,citecolor=BANavy,pdftitle={Six Knowledge Areas and 30 Tasks}}
\setlength{\parindent}{0pt}
\setlength{\parskip}{4pt}
\setlist{nosep,leftmargin=*}
\renewcommand{\arraystretch}{1.18}
\newcolumntype{Y}{>{\RaggedRight\arraybackslash}X}
\newcolumntype{L}[1]{>{\RaggedRight\arraybackslash}p{#1}}
\pagestyle{fancy}
\setlength{\headheight}{14pt}
\fancyhf{}
\lhead{\textcolor{BANavy}{\small BABOK v3 Study Series}}
\rhead{\textcolor{BASlate}{\small Six Knowledge Areas and 30 Tasks}}
\cfoot{\small \thepage\ / \pageref{LastPage}}
\renewcommand{\headrulewidth}{0.4pt}
\titleformat{\section}{\Large\bfseries\color{BANavy}}{\thesection}{0.7em}{}
\titleformat{\subsection}{\large\bfseries\color{BAOrange}}{\thesubsection}{0.6em}{}
\titleformat{\subsubsection}{\normalsize\bfseries\color{BANavy}}{\thesubsubsection}{0.5em}{}
\newtcolorbox{keybox}[1][]{enhanced,breakable,colback=BALight,colframe=BANavy,boxrule=0.7pt,arc=1.5mm,left=2mm,right=2mm,top=1.5mm,bottom=1.5mm,#1}
\newtcolorbox{practicebox}[1][]{enhanced,breakable,colback=white,colframe=BAOrange,boxrule=0.8pt,arc=1.5mm,left=2mm,right=2mm,top=1.5mm,bottom=1.5mm,#1}
\newtcolorbox{warningbox}[1][]{enhanced,breakable,colback=white,colframe=BAGold,boxrule=0.8pt,arc=1.5mm,left=2mm,right=2mm,top=1.5mm,bottom=1.5mm,#1}
\begin{document}
\begin{titlepage}
\centering
\vspace*{1.15in}
{\Huge\bfseries\color{BANavy} Six Knowledge Areas and 30 Tasks\par}
\vspace{0.35in}
{\Large\color{BAOrange} Purpose, outputs, and working questions\par}
\vspace{0.55in}
\rule{0.78\textwidth}{1.2pt}\par
\vspace{0.35in}
{\large Technical reference and study artifact based on the uploaded \emph{A Guide to the Business Analysis Body of Knowledge (BABOK Guide), Version 3}.\par}
\vfill
\begin{keybox}
\textbf{Source discipline.} BABOK terminology and organization are preserved where practical. Explanations are paraphrased for study and engineering use. Any engineering, security, Agile, UML/BPMN, or enterprise-architecture crosswalks are practitioner mappings, not normative BABOK statements.
\end{keybox}
\vspace{0.35in}
{\small Prepared October 2026\par}
\end{titlepage}
\tableofcontents
\newpage
\section{Quick Reference}
The knowledge areas are logical domains of expertise, not chronological project phases. The 30 tasks may be reordered, repeated, or performed concurrently according to the initiative and available inputs.
\begin{longtable}{L{0.28\textwidth} L{0.55\textwidth}}
\toprule
\textbf{Knowledge area} & \textbf{Primary concern}\\
\midrule
Business Analysis Planning and Monitoring & Organize and coordinate business analysis work, stakeholder participation, governance, information management, and improvement.\\
Elicitation and Collaboration & Obtain information, confirm what was learned, communicate it effectively, and sustain stakeholder collaboration throughout the work.\\
Requirements Life Cycle Management & Manage requirements and designs from inception through retirement, preserving alignment, traceability, priority, change control, and agreement.\\
Strategy Analysis & Understand the business need and current state, define the desired future state, assess uncertainty, and select a change strategy.\\
Requirements Analysis and Design Definition & Transform elicited information into coherent requirements and designs, assess quality and value, organize the architecture, and recommend a solution.\\
Solution Evaluation & Measure how an implemented or partially implemented solution performs, diagnose limitations, and recommend actions that increase realized value.\\
\bottomrule
\end{longtable}
\section{Business Analysis Planning and Monitoring}
Organize and coordinate business analysis work, stakeholder participation, governance, information management, and improvement.
\begin{longtable}{L{0.25\textwidth} L{0.39\textwidth} L{0.27\textwidth}}
\toprule
\textbf{Task} & \textbf{Intent} & \textbf{Typical information/artifact focus}\\
\midrule
\endhead
3.1 Plan Business Analysis Approach & Define how business analysis will be performed for the initiative. & BA approach, method, timing, deliverables, level of formality\\
3.2 Plan Stakeholder Engagement & Plan how to establish and maintain productive stakeholder relationships. & Stakeholder analysis, engagement strategy, communication and collaboration needs\\
3.3 Plan Business Analysis Governance & Define how requirements and design decisions, changes, priorities, reviews, and approvals will be governed. & Decision rights, review and approval paths, change control, prioritization rules\\
3.4 Plan Business Analysis Information Management & Define how business analysis information will be captured, organized, stored, accessed, traced, and retained. & Information structure, repositories, access, versioning, traceability approach\\
3.5 Identify Business Analysis Performance Improvements & Assess the effectiveness of business analysis work and identify improvements. & Performance measures, lessons, corrective actions, continuous improvement\\
\bottomrule
\end{longtable}
\begin{practicebox}\textbf{Working questions.}\begin{itemize}
\item How will analysis be performed and tailored?
\item Who must participate and how?
\item Who can decide, approve, or authorize change?
\item Where will analysis information live and how will it be traced?
\item How will the analysis process itself improve?
\end{itemize}\end{practicebox}
\section{Elicitation and Collaboration}
Obtain information, confirm what was learned, communicate it effectively, and sustain stakeholder collaboration throughout the work.
\begin{longtable}{L{0.25\textwidth} L{0.39\textwidth} L{0.27\textwidth}}
\toprule
\textbf{Task} & \textbf{Intent} & \textbf{Typical information/artifact focus}\\
\midrule
\endhead
4.1 Prepare for Elicitation & Set the scope of an elicitation activity, select techniques, and arrange the people, materials, and logistics needed. & Elicitation plan, objectives, participants, questions, supporting material\\
4.2 Conduct Elicitation & Explore and draw out information relevant to the change. & Raw elicitation results, observations, decisions, assumptions, open questions\\
4.3 Confirm Elicitation Results & Check elicited information for accuracy and consistency before relying on it. & Confirmed notes, reconciled findings, identified gaps or conflicts\\
4.4 Communicate Business Analysis Information & Present business analysis information so stakeholders develop a shared understanding. & Audience-appropriate models, documents, briefings, decisions\\
4.5 Manage Stakeholder Collaboration & Sustain stakeholder participation and alignment toward common outcomes. & Resolved issues, engagement actions, collaboration agreements\\
\bottomrule
\end{longtable}
\begin{practicebox}\textbf{Working questions.}\begin{itemize}
\item What information is missing or uncertain?
\item Which elicitation technique best fits the information need?
\item How will results be confirmed?
\item How should the information be communicated to each audience?
\item What threatens effective stakeholder collaboration?
\end{itemize}\end{practicebox}
\section{Requirements Life Cycle Management}
Manage requirements and designs from inception through retirement, preserving alignment, traceability, priority, change control, and agreement.
\begin{longtable}{L{0.25\textwidth} L{0.39\textwidth} L{0.27\textwidth}}
\toprule
\textbf{Task} & \textbf{Intent} & \textbf{Typical information/artifact focus}\\
\midrule
\endhead
5.1 Trace Requirements & Maintain relationships among needs, requirements, designs, solution components, and related work so changes and coverage can be assessed. & Bidirectional trace links, coverage and allocation views, impact paths\\
5.2 Maintain Requirements & Keep requirements and designs accurate, current, consistent, and reusable across their life cycle. & Current baselines, maintained requirements repository, reusable assets\\
5.3 Prioritize Requirements & Order requirements and designs according to relative importance, value, urgency, risk, and dependencies. & Priority decisions, sequencing, backlog order\\
5.4 Assess Requirements Changes & Evaluate the consequences of proposed changes to requirements and designs. & Impact assessment, recommended disposition, affected links and stakeholders\\
5.5 Approve Requirements & Obtain the level of agreement and authorization needed to continue analysis or proceed with solution work. & Approval record, decision, agreed baseline\\
\bottomrule
\end{longtable}
\begin{practicebox}\textbf{Working questions.}\begin{itemize}
\item What must be traceable to what?
\item Which requirements or designs need maintenance or reuse?
\item What is most important now and why?
\item What changes if this requirement changes?
\item What level of approval is actually required?
\end{itemize}\end{practicebox}
\section{Strategy Analysis}
Understand the business need and current state, define the desired future state, assess uncertainty, and select a change strategy.
\begin{longtable}{L{0.25\textwidth} L{0.39\textwidth} L{0.27\textwidth}}
\toprule
\textbf{Task} & \textbf{Intent} & \textbf{Typical information/artifact focus}\\
\midrule
\endhead
6.1 Analyze Current State & Establish why change is needed and understand the existing enterprise conditions affected by it. & Current-state capabilities, processes, constraints, problems, opportunities\\
6.2 Define Future State & Describe the conditions, goals, objectives, and capabilities that would satisfy the business need. & Future-state description, measurable objectives, target capabilities\\
6.3 Assess Risks & Evaluate uncertainty that could affect value during transition or in the future state, and identify suitable responses. & Risk assessment, assumptions, response recommendations\\
6.4 Define Change Strategy & Compare ways to move from current to future state and recommend the most valuable feasible approach, including transition states. & Gap analysis, option assessment, transition roadmap, recommended strategy\\
\bottomrule
\end{longtable}
\begin{practicebox}\textbf{Working questions.}\begin{itemize}
\item Why is change needed?
\item What does success look like?
\item What uncertainty could prevent value?
\item What path best moves the enterprise from current to future state?
\end{itemize}\end{practicebox}
\section{Requirements Analysis and Design Definition}
Transform elicited information into coherent requirements and designs, assess quality and value, organize the architecture, and recommend a solution.
\begin{longtable}{L{0.25\textwidth} L{0.39\textwidth} L{0.27\textwidth}}
\toprule
\textbf{Task} & \textbf{Intent} & \textbf{Typical information/artifact focus}\\
\midrule
\endhead
7.1 Specify and Model Requirements & Analyze, synthesize, and refine elicitation results into useful requirements and design representations. & Structured requirements, models, business rules, scenarios, data/process views\\
7.2 Verify Requirements & Check that requirements and designs meet quality expectations and are fit for their intended analytical or delivery purpose. & Quality review findings, corrected specifications and models\\
7.3 Validate Requirements & Confirm that requirements and designs support business requirements and the delivery of needed value. & Validated requirements/designs, rationale, value alignment\\
7.4 Define Requirements Architecture & Organize requirements so their relationships and collective structure support the objectives. & Requirements architecture, relationships, views, dependencies\\
7.5 Define Design Options & Develop solution approaches, allocate requirements, and represent viable options for achieving the future state. & Design options, allocation decisions, solution components\\
7.6 Analyze Potential Value and Recommend Solution & Estimate the potential value of design options and recommend the option best aligned with enterprise needs. & Option comparison, value assessment, solution recommendation\\
\bottomrule
\end{longtable}
\begin{practicebox}\textbf{Working questions.}\begin{itemize}
\item What representation best expresses the requirement or design?
\item Is the specification internally high quality?
\item Does it support the business need and value?
\item How do requirements relate and depend on one another?
\item Which design options are viable?
\item Which option provides the best potential value?
\end{itemize}\end{practicebox}
\section{Solution Evaluation}
Measure how an implemented or partially implemented solution performs, diagnose limitations, and recommend actions that increase realized value.
\begin{longtable}{L{0.25\textwidth} L{0.39\textwidth} L{0.27\textwidth}}
\toprule
\textbf{Task} & \textbf{Intent} & \textbf{Typical information/artifact focus}\\
\midrule
\endhead
8.1 Measure Solution Performance & Define meaningful measures and collect data to evaluate how effectively the solution delivers value. & Performance measures, observed data, measurement results\\
8.2 Analyze Performance Measures & Interpret performance information to understand realized value and whether current business needs are being met. & Performance insights, gaps, trends, conclusions\\
8.3 Assess Solution Limitations & Identify internal solution factors that constrain value realization. & Solution limitation analysis, defect/design/process constraints\\
8.4 Assess Enterprise Limitations & Identify external enterprise factors that prevent full value realization. & Enterprise constraint analysis, organizational/process/capability barriers\\
8.5 Recommend Actions to Increase Solution Value & Recommend actions that close the gap between potential and actual value. & Improvement recommendations, replacement/retirement/optimization options\\
\bottomrule
\end{longtable}
\begin{practicebox}\textbf{Working questions.}\begin{itemize}
\item What evidence indicates solution performance?
\item What does the evidence mean for realized value?
\item Which limitations are inside the solution?
\item Which limitations are in the enterprise context?
\item What action would most improve value?
\end{itemize}\end{practicebox}
\section{Relationship Patterns Across Knowledge Areas}
\begin{itemize}
\item \textbf{Strategy $\leftrightarrow$ RADD:} current/future-state reasoning provides context for requirements and design; emerging design feasibility can in turn change the strategy.
\item \textbf{Elicitation $\leftrightarrow$ RADD:} elicitation supplies information; modelling exposes gaps that trigger more elicitation.
\item \textbf{RLCM $\leftrightarrow$ all areas:} traceability, prioritization, maintenance, change assessment, and approval govern evolving information across the initiative.
\item \textbf{Solution Evaluation $\leftrightarrow$ Strategy:} actual performance can reveal a new need, a changed context, or a need to revise the future state.
\item \textbf{Planning and Monitoring $\rightarrow$ all areas:} governance, information management, stakeholder strategy, and performance improvement influence how every other task is carried out.
\end{itemize}

\section*{Primary Source}
\addcontentsline{toc}{section}{Primary Source}
International Institute of Business Analysis (IIBA), \emph{A Guide to the Business Analysis Body of Knowledge (BABOK Guide)}, Version 3, 2015. This document is a study and practitioner synthesis of the user-supplied copy; it is not an official IIBA publication.
\end{document}
